\documentclass[runningheads]{llncs}

\usepackage[T1]{fontenc}
\usepackage{amsmath,amssymb}
\usepackage{graphicx}
\usepackage{booktabs}
\usepackage{xspace}
\usepackage[hidelinks]{hyperref}

% Notation shared by all sections.
\newcommand{\pstar}{p^{*}}
\newcommand{\Vmax}{V_{\max}}
\newcommand{\F}{\mathbb{F}}
\newcommand{\eg}{e.g.,\xspace}
\newcommand{\ie}{i.e.,\xspace}

\begin{document}

% TODO(author): confirm title and author list (taken from the FinteQC paper).
\title{Zeequent Batch Auctions}
\titlerunning{Zeequent Batch Auctions}

\author{Kimia Esmaili\inst{1} \and Elizabeth van Oorschot\inst{2} \and Jeremy Clark\inst{1}}
\authorrunning{K. Esmaili, E. van Oorschot, and J. Clark}
\institute{Concordia University, Montreal, Canada \and McGill University, Montreal, Canada}

\maketitle

% The abstract is the author's own text and is kept as given.
% TODO(author): "100 to 1,000 price ticks" (results cover 100 to 8,000 in Rust,
% 100 to 2,000 in Noir) and "arcworks" (arkworks).
\begin{abstract}
Frequent batch auctions (FBAs) are an alternative market structure to continuous limit order books (CLOBs) in which, instead of immediate trade matching, periodic uniform-price clearing occurs over short time intervals. FBAs reduce the value of sub-second latency advantages, which are exploited by high-frequency traders (HFTs). The issue with this market structure is that participants never see the underlying data used to compute the final clearing price; hence, its soundness is questionable. In this paper, we propose a zero-knowledge FBA protocol and provide a construction in which the person announcing the final clearing price at an exchange proves with a custom zk-SNARK---without revealing individual orders---that a published clearing price and allocation satisfy the auction rule. Thus, this protocol preserves participant privacy while providing public verifiability for the end result. We arithmetize the FBA clearing process as vanishing polynomial identities over a multiplicative subgroup, built from reusable gadgets and verified via KZG10 commitments with Fiat-Shamir batching. Beyond theoretical formulation, we implemented the protocol in both Rust/arcworks and Noir/Barretenberg, each proved at 100 to 1,000 price ticks, using both our own synthetic data and real data from Nasdaq TotalView, demonstrating its practical utility.

\keywords{frequent batch auctions \and call markets \and zk-SNARKs \and polynomial IOPs}
\end{abstract}

\section{Introduction}
\label{sec:intro}

Modern stock exchanges such as Nasdaq and NYSE run two trading systems~\cite{harris2003}. During regular hours, traders submit orders to a continuous limit order book (CLOB), where a buy order and a sell order execute as soon as their prices meet. Traders can also submit orders to a scheduled auction, or call market, which holds orders until a fixed time and then executes them together at a single price. The closing auction sets the official closing price of each stock, against which portfolios are benchmarked, indexes are computed, and derivatives are settled; closing auctions account for 7.5\% of daily volume in U.S.\ stocks~\cite{bogousslavsky2023}. The same mechanism opens the market the next morning.

Most activity still occurs in the CLOB, and a CLOB processes orders in the order they arrive. A difference of microseconds therefore decides who trades first. High-frequency traders (HFTs) invest in fast hardware, private data links, and co-location to win these races. Aquilina, Budish, and O'Neill measure them directly on exchange message data: races account for about 20\% of trading volume in FTSE~100 stocks and the sums at stake are on the order of \$5~billion a year across global equity markets~\cite{aquilina2022}. HFTs also narrow spreads and add liquidity, but liquidity that can be withdrawn faster than other traders can react is of limited use to them.

\paragraph{Frequent batch auctions.}
In an opening or closing auction, reaching the exchange a few microseconds earlier brings no advantage, because orders are not matched one by one. Budish, Cramton, and Shim propose to run such an auction repeatedly through the day~\cite{budish2015,budish2014}. In a frequent batch auction (FBA), the exchange collects orders over a short interval (\eg one second) and then clears the whole batch at one price. An order that arrives at the start of the batch is treated the same as one that arrives at the end: there is price priority and no time priority. Traders compete on price and no longer on speed.

\paragraph{The transparency gap.}
An FBA gives up something a CLOB has. In a CLOB the book is visible, so once an order reaches the matching engine there is no doubt about why it executed, did not execute, or executed at a given price. In a batch auction, no participant sees the set of orders that determines the clearing price, and after the close participants see only the outcome. The orders that failed to execute stay hidden, although the clearing price depended on them. A trader observes whether their own order executed and at what price, and cannot check that the exchange included every eligible order, excluded ineligible ones, chose the price the rule prescribes, and allocated trades correctly. The exchange cannot repair this by publishing the book, since unexecuted orders reveal trading interest that their owners have reason to keep private.

\paragraph{A proof instead of the book.}
A zero-knowledge proof~\cite{gmw1986} shows that a computation was performed correctly without revealing its private inputs. Applied to a batch auction, the exchange proves that the published clearing price and volumes follow from the submitted orders and reveals none of them. Custodial cryptocurrency exchanges already publish proofs of this kind about their hidden balance sheets~\cite{provisions2015}; we apply the approach to the matching engine.

We develop a zero-knowledge frequent batch auction (zk-FBA). The exchange commits to the order book as orders arrive, returns to each trader a receipt that the order was recorded, and publishes at the close a clearing receipt with a succinct proof (a zk-SNARK) that the receipt is the outcome the auction rule prescribes for the committed book. The exchange must produce this proof once per batch, which is the binding constraint on the design. A general-purpose proof system compiles the clearing rule into a circuit and pays for the generality at every batch. We instead write the rule directly as polynomial identities over tick-indexed columns: a custom arithmetization~\cite{plonkbook} for a polynomial interactive oracle proof (PIOP)~\cite{dark2020}, in the style of PLONK~\cite{plonk2019} and with KZG commitments~\cite{kzg2010}.

An FBA has no canonical specification, since it is a market-design proposal and not an operational standard~\cite{budish2014}. We fix the details by adopting the rule of a deployed auction, the Nasdaq Closing Cross, and assume only that it runs frequently.

\subsection{Contributions}
\label{sec:contributions}

\begin{enumerate}
  \item \textbf{A custom arithmetization of the clearing rule.} We encode a closing-cross clearing rule, with limit, market, and imbalance-only orders and its three tie-breaks, as 33 polynomial constraints over 88 columns indexed by price tick, combined into one quotient. We prove that the constraints admit exactly one clearing receipt for a given book (\S\ref{sec:protocol}).
  \item \textbf{Lookup range checks.} Every inequality in the rule is a range check. We replace one committed column per bit with 8-bit limbs and a single logarithmic-derivative lookup, which reduces the proof for 32-bit volumes from 352 commitments and 22.9\,KB to 89 commitments and 6.1\,KB (\S\ref{sec:protocol}).
  \item \textbf{A committed, updatable book.} We commit to the book with a homomorphic polynomial commitment, so that the exchange updates the commitment for each order, returns a receipt to the trader, and hides the side of the order in the public log. A second proof at the close shows that every logged update is a valid order (\S\ref{sec:book}).
  \item \textbf{Two implementations.} We implement the prover and verifier in Rust on arkworks, and independently write the clearing statement as a Noir circuit proved with Barretenberg. For a book of 1,000 price ticks, we generate the clearing proof in 1.34\,s on one core of a laptop and 0.37\,s on eight threads; the proof is 6,124 bytes and verifies in 29\,ms. This is 5.5 times faster than the circuit, with a proof 2.4 times smaller (\S\ref{sec:performance}).
  \item \textbf{Order flow from a real market.} No batch-auction venue publishes order-level data. We replay Nasdaq TotalView-ITCH order flow as one-second batches, measure how often each tie-break is reached, and build a calibrated generator of benchmark books. Both implementations give the same receipt on all 32 real and 25 synthetic books (\S\ref{sec:performance}).
\end{enumerate}

\section{Background and Related Work}
\label{sec:background}

\subsection{Batch Auctions and the Closing Cross}
\label{sec:fba}

We take the clearing rule from the Nasdaq Closing Cross~\cite{nasdaq4754} and assume it runs once per batch. Each asset (an equity such as AAPL, a token such as ETH, a futures contract) has its own auction. Prices lie on a grid of $N$ ticks.

\paragraph{Order types.}
We support three order types, each as a buy or a sell.\footnote{These correspond to limit-on-close (LOC), market-on-close (MOC), and imbalance-only (IO) orders in the Closing Cross.} A \emph{limit order} carries a limit price $\ell$. With $\pstar$ the clearing price of the batch, a buy limit order (bid) is eligible to execute only if $\pstar \le \ell$, and a sell limit order (ask) only if $\pstar \ge \ell$. A \emph{market order} carries only a quantity and is eligible at any $\pstar$. An \emph{imbalance-only (IO) order} has a limit price like a limit order, but is eligible only when it reduces the residual imbalance at $\pstar$: a buy IO requires $\pstar \le \ell$ and excess supply at $\pstar$, and a sell IO requires $\pstar \ge \ell$ and excess demand.

\paragraph{Clearing.}
At the end of a batch the exchange determines $\pstar$ in three phases.

\emph{Phase 1.} Using limit and market orders only, the exchange computes for each tick $p$ the cumulative buy depth $D(p)$, the cumulative sell depth $S(p)$, and the executable volume $V(p) = \min(D(p), S(p))$. Let $\Vmax = \max_p V(p)$. Since $D$ is non-increasing and $S$ is non-decreasing, $V$ rises weakly to the crossing region and then falls weakly, so the maximizing ticks form a contiguous interval $[c,d] = \{p : V(p) = \Vmax\}$, which we call the \emph{plateau}. If $c = d$, this tick is $\pstar$.

\emph{Phase 2.} Otherwise the exchange computes the imbalance $\Delta(p) = |D(p) - S(p)|$ on the plateau and keeps the ticks that minimize it. If one tick remains, it is $\pstar$; a deterministic tertiary rule decides among several (\S\ref{sec:protocol}).

\emph{Phase 3.} The exchange executes all eligible limit and market orders at $\pstar$. Orders priced strictly better than $\pstar$ execute in full, and such limit orders receive price improvement. If buy and sell volume at $\pstar$ are unequal, one side is long and the other short. The exchange admits eligible IO orders on the short side to reduce the imbalance, and rations what remains of the long side's interest at exactly $\pstar$ pro rata.

The outcome of a batch is a \emph{clearing receipt} of eight public values (\S\ref{sec:protocol}), or the statement that no bid meets an ask.

\subsection{Succinct Proofs}
\label{sec:snarks}

Zero-knowledge proofs date to the 1980s~\cite{gmw1986,bcc1988}, and classic constructions produce proofs at least linear in the size of the statement. A succinct argument (SNARK) is sublinear, typically of logarithmic or constant size. Several SNARKs~\cite{sonic2019,marlin2020,plonk2019} share the framework of a polynomial interactive oracle proof~\cite{dark2020}: the prover specifies polynomials, and the verifier checks identities among them at random points. A polynomial commitment scheme makes this concrete by binding the prover to each polynomial and proving its evaluations. We use KZG~\cite{kzg2010}, whose commitments and opening proofs are single group elements. KZG needs a structured reference string $(g, g^{\tau}, g^{\tau^2}, \dots)$ for a secret $\tau$ that must be discarded; the string is universal, and public ceremonies in which one honest participant suffices have produced it at scale~\cite{nikolaenko2024}. The Fiat--Shamir transform~\cite{fiatshamir1986} removes the verifier's messages, provided the transcript binds the full statement~\cite{bernhard2012,dao2023}.

\paragraph{Custom arguments.}
A computation reaches a SNARK in one of three ways: as a program compiled to the instructions of a virtual machine, as a circuit written in a domain-specific language, or as an arithmetization designed for that computation alone. Each level of abstraction costs the prover time. For ballot validity, a tailored proof is expected to be faster than a general-purpose one~\cite{huber2024}, and Zeeperio's custom PIOP for election tallies proves orders of magnitude more ballots than circuit-based systems in the literature~\cite{zeeperio2026}. We follow the third route. A batch auction suits it: the state is a small number of arrays indexed by price tick, and the rule is a fixed sequence of prefix sums, minima, comparisons, and selections, each of which is a known PIOP gadget~\cite{plonkbook}.


\subsection{Related Work}
\label{sec:related}

\paragraph{Secrecy-preserving auctions and exchanges.}
Parkes et al.\ run sealed-bid auctions in which the auctioneer proves the outcome correct over encrypted bids~\cite{parkes2008}, and Thorpe and Parkes extend the idea to a securities exchange whose operator proves facts about a hidden order book, such as the spread or the depth~\cite{thorpe2007}. Their proofs grow with the number of orders. A second line removes the operator. Bogetoft et al.\ compute the clearing price of a double auction among three parties with secure multiparty computation, in a deployment for Danish sugar-beet contracts~\cite{bogetoft2009}. Cartlidge, Smart, and Talibi Alaoui evaluate multiparty computation for dark pools and find volume matching practical where a continuous double auction is not~\cite{cartlidge2019}; FuturesMEX distributes a futures exchange among its traders~\cite{massacci2018}, and Ngo et al.\ build dark-pool trading on a witness-key-agreement primitive~\cite{ngo2021}. A recent systematization covers 22 such protocol families~\cite{arote2026}.

\paragraph{Periodic and batch auctions.}
Closest to our setting, Galal and Youssef give a publicly verifiable periodic auction that keeps bids secret, using aggregated Bulletproofs~\cite{bulletproofs2018} to prove the ordering of a vector of commitments~\cite{galal2021}. FairTraDEX builds a decentralised exchange on frequent batch auctions from zero-knowledge set-membership proofs and an escrowed commit--reveal step, with game-theoretic guarantees against value extraction~\cite{fairtradex2022}; orders are revealed when the batch ends. We differ in three respects. We keep the operator, as venues in regulated markets do, and make its output verifiable. We prove the full rule of a deployed cross, including imbalance-only orders and the tie-breaks that select one price from a plateau. Our proofs have constant size in the number of orders and price ticks. We are not aware of prior work with these three properties together.

\paragraph{Proofs by exchanges.}
Proofs of solvency are the established case of an exchange proving a statement about data it keeps private. Provisions gave the first zero-knowledge construction~\cite{provisions2015}; later systems shorten the proofs, and an examination of deployed ones found several broken~\cite{chalkias2022}. Xiezhi builds an end-to-end proof of solvency as a custom polynomial argument over KZG~\cite{xiezhi2026}. Our protocol applies the same method to a different duty of the exchange, and the two compose: one venue can prove both that it clears correctly and that it holds what it owes.

% \input{sections/03-protocol}
% \input{sections/04-book}
% \input{sections/05-performance}
% \input{sections/06-security}
% \input{sections/07-discussion}

\bibliographystyle{splncs04}
\bibliography{refs}

% \appendix
% \input{sections/A-clearing-rule}

\end{document}
